\documentclass[12pt,letterpaper, onecolumn]{exam}
\usepackage{amsmath}
\usepackage{amssymb}
\usepackage{amsthm}
\usepackage{graphicx}
\usepackage{mathrsfs}
\usepackage{xcolor}
\usepackage[lmargin=71pt, tmargin=1.2in]{geometry}  %For centering solution box
\usepackage[english]{babel}
\newtheorem{theorem}{Theorem}
\newtheorem{lemma}[theorem]{Lemma}
\usepackage{accents}
\usepackage{ulem} % sout() to crossing things out 
\newcommand{\ubar}[1]{\underaccent{\bar}{#1}}
\newcommand{\bs}[1]{\boldsymbol{#1}}
\newcommand{\bm}[1]{\mathbf{#1}}
\newcommand{\der}[2]{\frac{d #1}{d#2}}
\newcommand{\pder}[3]{\frac{\delta^{#3} #1}{\delta^{#3}#2}}
\newcommand{\rs}{$X_1, ..., X_n$}
\newcommand{\nsum}{\sum_{i=1}^n}
\newcommand{\nprod}{\prod_{i=1}^n}
\newcommand*{\Comb}[2]{{}^{#1}C_{#2}}%
\newcommand{\ifc}{\textrm{ if }}
\newcommand{\real}{\mathbb{R}}
\newcommand{\nat}{\mathbb{N}}
\newcommand{\integer}{\mathbb{Z}}
\newcommand{\power}{\mathcal{P}}
\newcommand{\rational}{\mathbb{Q}}
\newcommand{\cutelim}{\lim\limits_{n\to\infty}}
\newcommand{\salg}{$\sigma$-algebra }
\newcommand{\I}{\mathbf{I}}

\newcommand{\calA}{\mathcal{A}}
\newcommand{\calO}{\mathcal{O}}
\newcommand{\calC}{\mathcal{C}}
\newcommand{\calB}{\mathcal{B}}
\newcommand{\calI}{\mathcal{I}}
\newcommand{\calF}{\mathcal{F}}
\newcommand{\calG}{\mathcal{G}}
\newcommand{\gensig}[1]{\sigma\langle #1 \rangle} % notation for a sigma algebra generated by some set 
\newcommand{\om}{\mu^*} % outer measure 
\newcommand{\rat}{\mathbb{Q}}
\newcommand{\tribrac}[1]{\langle #1 \rangle}
\newcommand{\borel}{\mathcal{B}(\mathbb{R})}
\newcommand{\boreltwo}{\mathcal{B}(\mathbb{R}^2)}
\newcommand{\borelk}[1]{\mathcal{B}(\mathbb{R}^#1)}
\newcommand{\probspace}{(\Omega, \calF, P)}

\newcommand{\Iunion}{\bigcup_{\alpha \in I}}
\newcommand{\Iintersect}{\bigcap_{\alpha \in I}}

\newcommand{\dm}{d \mu}

\newcommand{\cutesum}[1]{\sum_{#1 = 1}^\infty}


% \chead{\hline} % Un-comment to draw line below header
\thispagestyle{empty}   %For removing header/footer from page 1

\begin{document}

\begingroup  
    \centering
    \LARGE STAT 6410\\
    \LARGE Homework 8\\[0.5em]
    \large Sarah Kilgard\par
\endgroup
\rule{\textwidth}{0.4pt}
\pointsdroppedatright   %Self-explanatory
\printanswers
\renewcommand{\solutiontitle}{\noindent\textbf{Ans:}\enspace}   %Replace "Ans:" with starting keyword in solution box

\begin{questions}
%%%%%%%%%%% Question 1 %%%%%%%%%%%%%%%%%%%%    
    \question \textbf{Chapter 4, Problem 3}

    Let $\nu$ be the counting measure on $[0,1]$ and $\int_{[0,1]} h d\nu < \infty$ for some nonnegative function $h$. Show that $B = \{x : h(x) > 0\}$ is countable. 

    (\textbf{Hint:} Let $B_n = \{x : h(x) > \frac{1}{n}\}$. Show that $B_n$ is a finite set for each $n \in \nat$.)
    \begin{solution}
    Clearly, $B_n = \{x \in [0, 1]: h(x) > 0\} = \subseteq [0, 1]$ for all $n \in \nat$. \\ Thus $\int_{B_n} h d\nu \leq \int_{[0,1]} h d\nu < \infty$ for all $n \in \nat$. 

    If $x \in B_n$, $\frac{1}{n} < h(x)$, so $\int \frac{1}{n} I_{B_n} < \int h I_{B_n} d\nu = \int_{B_n} h d \nu \leq \int_{[0,1]} h d\nu < \infty$.

    $\int \frac{1}{n} I_{B_n} d\nu = \frac{1}{n} \nu(B_n) < \infty \implies \nu(B_n) < \infty$  for all $n \in \nat$.

    Since $B_n$ is finite for every $n \in \nat$ and $\cup_{n\geq 1} B_n = B$, $B$ is countable. 
    \end{solution}

%%%%%%%%%%% Question 2 %%%%%%%%%%%%%%%%%%%%    
    \question \textbf{Chapter 4, Problem 4}

    Prove Proposition 4.1.2.

    \textbf{Proposition 4.1.2}: Let $\nu$. $\mu$, $\mu_1$, $\mu_2$, ... be $\sigma$-finite measure on a measurable space $(\Omega, \calF)$
    \begin{itemize}
        \item[(i)] If $\mu_1 << \mu_2$ and $\mu_2 << \mu_3$, then $\mu_1 << \mu_3$ and
        \[\frac{d \mu_1}{d \mu_3} = \frac{d \mu_1}{d \mu_2}\frac{d \mu_2}{d \mu_3} \: a.e.(\mu_3)\]
        \item[(ii)] Suppose that $\mu_1$ and $\mu_2$ are dominated by $\mu_3$. Then for any $\alpha, \beta \geq 0$, $\alpha \mu_1 + \beta \mu_2$ is dominated by $\mu_3$ and 
        \[\frac{d(\alpha \mu_1 + \beta \mu_2)}{d \mu_3} = \alpha \frac{d \mu_1}{d \mu_3} + \beta \frac{d \mu_2}{d \mu_3} a.e.(\mu_3)\]
        \item[(iii)] If $\mu << \nu$ and $\frac{d \mu}{d \nu} > 0 \: a.e.(\nu)$, then $\nu << \mu$ and 
        \[\frac{d \nu}{d \mu} = \left(\frac{d \mu}{d \nu}\right)^{-1} \: a.e.(\mu).\]
        \item[(iv)] Let $\{\mu_n\}_{n \geq 1}$ be a sequence of measure and $\{\alpha\}_{n \geq 1}$ be a sequence of positive real numbers, i.e., $\alpha > 0$ for all $n \geq 1$. Define $\mu = \cutesum{n} \alpha_n \mu_n$
        \begin{itemize}
            \item[(a)] Then $\mu << \nu$ iff $\mu_n << \nu$ for each $n \geq 1$ and in this case, 
            \[\frac{d \mu}{d \nu} = \cutesum{n}\alpha_n \frac{d \mu_n}{d \nu} \: a.e.(\nu).\]
            \item[(b)] $\mu \perp \nu$ iff $\mu_n \perp \nu$ for all $n \geq 1$ 
        \end{itemize}
    \end{itemize}
    
    \begin{solution}
    \textbf{Lemma:} For any $\sigma$-finite measures $\mu, \: \nu$ on $(\Omega, \calF)$ with $\mu << \nu$, for any measurable function h: 
            $$\int h d\mu = \int h \frac{d\mu}{d\nu}d\nu$$
            \begin{proof}
                \underline{1: Indicator Functions} Let $h = \I_{A}$ for some $A \in \calF$. \\
                \begin{align*}
                    & \int h d\mu = \int \I_A d\mu = \int_A d\mu \\
                    & = \int \frac{d\mu}{d\nu}d\nu && \textrm{(by RN Theorem, where $\tfrac{d \mu}{d \nu}$ is the RN der.)}  \\
                    & = \int \I_A \frac{d\mu}{d\nu}d\mu = \int h \frac{d\mu}{d\nu}d\nu
                \end{align*}
                \underline{2: Simple Functions} Let $\{A_1, A_2, ... , A_k\} \subset \calF$ and $c_1, c_2, ... , c_k \in \real$. Define $h = \sum_{i = 1}^k c_i \I_{A_i}$
                \begin{align*}
                    & \int h d\mu = \int \sum_{i = 1}^k c_i\I_{A_i} d\mu = \sum_{i = 1}^k c_i \int \I_{A_i} d\mu \\
                    & = \sum_{i = 1}^k c_i \int \I_{A_i} \frac{d\mu}{d\nu}d\nu && \textrm{(\underline{1})}  \\
                    & = \int \sum_{i = 1}^k c_i \I_{A_i} \frac{d\mu}{d\nu}d\nu = \int h \frac{d\mu}{d\nu}d\nu
                \end{align*}
                \underline{3: Nonnegative Measurable Function, $h$, on $(\Omega, \calF)$:} Define $\{h_n\}_{\{n \geq 1\}}$ is a sequence of nonnegative simple functions such that $h_n \uparrow h$. 
                \begin{align*}
                    & \int h d\mu = \lim_{n \to \infty} \int h_n d\mu && (\textrm{def of int. of nonneg. mbl. fn}) \\
                    & = \lim_{n \to \infty} \int h_n \frac{d \mu}{d \nu}d\nu && (\underline{2}) \\
                    & = \int h \frac{d \mu}{d \nu}d\nu && (\textrm{MCT})
                \end{align*}
                \underline{4: Measurable Function on $(\Omega, \calF$), $h$} $h$ can be written as a linear combination of 2 nonnegative measurable functions, $h^+$ and $h^-$. $h = h^+ - h^-$. \\
                From (\underline{3}), we know that \\
                $\int h d \mu = \int h^+ d\mu - \int h^- d\mu = \int h^+ \frac{d \mu}{d \nu} d\nu - \int h^- \frac{d \mu}{d \nu} d\nu = \int h^+ - h^- \frac{d \mu}{d \nu} d\nu = \int h \frac{d \mu}{d \nu} d\nu$
            \end{proof}
        (i) \begin{proof}
            Suppose, without loss of generality, that $\mu_3(A) = 0$ for some $A \in \calF$. Thus, $\mu_2(A) = 0$, since $\mu_2 << \mu_3$. Finally, $\mu_1(A) = 0$, since $\mu_1 << \mu_2$. Therefore, $\mu_1 << \mu_3$. \\

            \begin{align*}
            & \mu_1(A) = \int I_A \mu_1 = \int I_A \frac{d \mu_1}{d \mu_2}d\mu_2 && \textrm{(R-N Theorem)} \\
            & \textrm{From R-N theorem, we know that } I_A \frac{d \mu_1}{d \mu_2} \textrm{ is a mble nonneg fn} \\
            & \mu_1(A) = \int I_A \frac{d \mu_1}{d \mu_2}d\mu_2 = \int I_A \frac{d \mu_1}{d \mu_2}\frac{d\mu_2}{d\mu_3}d\mu_3 && \textrm{(from the Lemma)} \\
            & \frac{d \mu_1}{d \mu_2}\frac{d\mu_2}{d\mu_3} \textrm{ is the R-N derivative of } \mu_1 \: w.r.t. \: mu_3 \\
            & \textrm{By R-N Theorem, } \mu_1(A) = \int I_A \frac{d \mu_1}{d \mu_3} d\mu_3 \\
            & \textrm{From the definition of an R-N derivative, we know: } \frac{d \mu_1}{d \mu_3} = \frac{d \mu_1}{d \mu_2}\frac{d\mu_2}{d\mu_3} a.e.(\mu_3)
            \end{align*}
        \end{proof}

        (ii) \begin{proof}
            Because $\alpha, \beta \geq 0$, $\alpha \mu_1 + \beta \mu_2$ is $\sigma$-finite. 

            For all $A \in \calF$
            \begin{align*}
                &\alpha\mu_1(A) + \beta \mu_2(A) = \alpha \int I_A d \mu_1 + \beta \int I_A d \mu_2 \\
                & = \alpha \int I_A \frac{d\mu_1}{d\mu_3} d\mu_3 + \beta \int I_A \frac{d\mu_2}{d\mu_3}d\mu_3 && \textrm{(R-N Theorem)} \\
                &  = \int_A \alpha \frac{d\mu_1}{d\mu_3}  +  \beta \frac{d\mu_2}{d\mu_3}d\mu_3
            \end{align*}
            Thus, $\alpha \frac{\mu_1}{\mu_3}  +  \beta \frac{\mu_2}{\mu_3}$ is the R-N derivative for $\alpha \mu_1 + \beta \mu_2$ w.rt. $\mu_3$

            Since $\mu_1 << \mu_3$ and $\mu_2 << \mu_3$, $\mu_3(A) = 0 \implies \mu_2(A) = 0, \: \mu_1(A) = 0$. Thus, $\alpha \mu_1(A) + \beta \mu_2(A) = 0$. Thus, $\alpha \mu_1 + \beta \mu_2 << \mu_3$.

            Thus, $\alpha\mu_1(A) + \beta \mu_2(A) = \int_A \frac{\alpha\mu_1 + \beta \mu_2}{d \mu_3} d \mu_3$ for all $A \in \calF$

            Thus, $\frac{d(\alpha\mu_1 + \beta \mu_2)}{d \mu_3}$ is the R-N derivative of $\alpha\mu_1 + \beta \mu_2$ w.r.t. $\mu_3$.

            By the definition of an R-N derivative, $\frac{d(\alpha\mu_1 + \beta \mu_2)}{d \mu_3} = \alpha \frac{\mu_1}{\mu_3}  +  \beta \frac{\mu_2}{\mu_3} a.e.(\mu_3)$.
            \end{proof}

            (iii) \begin{proof}
                
                Suppose for some $A \in \calF$ that $\mu(A) = 0$. \\
                Thus, $0 = \int I_A \frac{d \mu}{d \nu} d\nu$, since $\mu << \nu$. From R-N theorem, we know that $\frac{d \mu}{d \nu}$ is measurable, and from the problem, we know that $\frac{d \mu}{d \nu} > 0$. 

                Thus, $0 = \int I_A \frac{d \mu}{d \nu} d\nu = \int I_A d \nu = \nu(A)$. 

                Thus, $\nu << \mu$.

               For any $A \in \calF$, define, $h \equiv I_A \left(\frac{d \mu}{d\nu}\right)^{-1}$. We know $h$ is measurable and nonnegative from the R-N theorem ($\mu << \nu$) and because $\frac{d \nu}{d\nu} > 0 a.e. (\nu)$.

               \[\nu(A) = \int I_A d \nu = \int I_A \frac{d \nu}{d\mu} \frac{d\mu}{d\nu}d\nu =  \int I_A \left(\frac{d \mu}{d\nu}\right)^{-1} \frac{d\mu}{d\nu}d\nu \]\[= \int I_A \left(\frac{d \mu}{d\nu}\right)^{-1} d\mu \\
               = \int h d \mu = \int I_A \left(\frac{d \mu}{d\nu}\right)^{-1} d\mu\]

               \[\nu(A) = \int I_A \nu = \int I_A \frac{d \nu}{d\mu} d\mu\]
               
                Since $\frac{d \nu}{d\mu}$ and $\left(\frac{d \mu}{d\nu}\right)^{-1}$ are both R-N derivatives of $\nu$ with respect to $\mu$, $\frac{d \nu}{d\mu} = \left(\frac{d \mu}{d\nu}\right)^{-1} a.e.(\mu)$
            \end{proof}

            (iv) 

           $\mu$ is clearly a valid measure on $(\Omega, \calF)$. 

           (a)
           \begin{proof}
                Suppose $\mu << \nu$. Suppose, that for some arbitrary $A \in \calF$, $\nu(A) = 0$. Thus, $\mu(A) = 0 = \mu(\sum_{n = 1}^\infty \alpha_n \mu_n(A)$. Since, $\mu_n(A) \geq 0$ and $\alpha_n > 0$ for all $n \in \nat$, $\mu_n(A) = 0$ for all $n \geq 1$. Thus, $\mu_n << \nu$ for all $n \geq 1$.

                 Suppose $\mu_n << \nu$ for each $n \geq 1$. Suppose, that for some arbitrary $A \in \calF$, $\nu(A) = 0$. Thus, for each $n \geq 1$, $\mu_n(A) = 0$. $\mu(A) = \cutesum{n} \alpha_n \mu_n(A) = \cutesum{n} \alpha_n \times 0 = 0$. Thus, $\mu << \nu$.

                 \[\mu(A) = \int_A \frac{d \mu}{d\nu} d\nu\]
                 \[\mu(A) = \cutesum{n} \alpha \mu_n(A) = \cutesum{n} \int_A \alpha_n  \frac{d \mu_n}{d \nu} d\nu\]
                 From, R-N Theorem, we know that $\alpha_n \frac{d \mu_n}{d\nu}$ are nonnegative and measurable, so we can invoke MCT:
                 \[\mu(A) = \int_A \cutesum{n}\alpha_n \frac{d \mu_n}{d\nu} d\nu\]

                 By the definition of an R-N derivative, $\frac{d \mu}{d\nu} = \cutesum{n} \alpha_n \frac{d \mu_n}{d\nu} a.e.(\nu)$

                 Thus, $\mu << \nu \iff \mu_n << \nu$ for all $n \geq 1$
            \end{proof}

            (b) \begin{proof}
                Suppose that $\mu \perp \nu$. Thus, there exists some set $B \in \calF$ such that $\mu(B) = 0$ and $\nu(B^c) = 0$. As in (a), $\mu(B) = 0 \implies \mu_n(B) = 0$ for all $n \geq 1$.

                Thus, $\mu_n \perp \nu$ for all $n \geq 1$.

                Suppose that $\mu_n \perp \nu$ for all $n \geq 1$. Thus, there exists some collection of set $B_n \in \calF$ such that $\mu_n(B_n) = 0$ and $\nu(B_n^c) = 0$ for all $n \geq 1$. 
                
                $\cap_{n \geq 1} B_n \in \calF$. For each $n \geq 1$, $\cap_{n\geq1} B_n \subseteq B_n$. Thus, $\mu_n(B_n) = 0 \implies \mu_n(\cap_{n\geq 1}B_n) = 0$. $\mu(\cap_{n\geq 1} B_n) = \cutesum{n}\mu_n(\cap_{n\geq 1} B_n) = 0$

                $(\cap_{n\geq 1} B_n)^c = \cup_{n\geq1}B_n^c$ \\
                $\nu( \cup_{n\geq1}B_n^c) \leq \cutesum{n}\nu(B_n^c) = \cutesum{n} 0 = 0$ (Countable subadditivity)

                $\nu \perp \mu_n$ for all $n \geq 1$

                Thus, $\nu \perp \mu \iff \nu \perp \mu_n$ for all $n \geq 1.$
            \end{proof}
    \end{solution}

%%%%%%%%%%% Question 3 %%%%%%%%%%%%%%%%%%%%    
    \question \textbf{Chapter 4, Problem 5}
    Find the Lebesgue decomposition of $\mu$ w.r.t $\nu$ and the Radon-Nikodym derivative $\frac{d \mu_a}{d\nu}$ in the following cases where $\mu_a$ is the absolutely continous component of $\mu$ w.r.t. $\nu$.
    \begin{itemize}
        \item[(a)] $\mu = N(0, 1), \: \nu = Exponential(1)$
        \item[(b)] $\mu = Exponential(1)$, $\nu = N(0,1)$
        \item[(c)] $\mu = \mu_1 + \mu_2$, where $\mu_1 = N(0,1)$, $\mu_2 = Poisson(1)$ and $\nu = Cauchy(0, 1)$.
        \item[(d)] $\mu=\mu_1+\mu_2$, $\nu=Geometric(p)$, $0<p<1$,where $\mu_1 =N(0,1)$ and $\mu_2$ = Poisson(1).
        \item[(e)] $\mu=\mu_1+\mu_2$, $\nu=\nu_1+\nu_2$ where $\mu_1 =N(0,1)$, $\mu_2$ = Poisson(1), $\nu_1$ = Cauchy(0, 1) and $\nu_2$ = Geometric(p), $0 < p < 1.$
        \item[(f)] $\mu$ = Binomial (10, 1/2), $\nu$ = Poisson (1). The measures referred to above are defined in Tables 4.6.1 and 4.6.2, given at the end of this section. 
    \end{itemize}
    \begin{solution}
        (a)
        
        Define $m$ to be the Lebesgue measure on $(\real, \borel)$, and $m_1(A) \equiv m(A \cap [0, \infty))$ for all $A \in \borel$, i.e. the Lebesgue measure restricted to positive numbers. 

        Define $\mu_a(A) \equiv \mu(A \cap [0, \infty))$ for all $A \in \borel$. In a similar manner to 4.1.1, we can show that $\mu_a << m_1$ and $m_1 << \nu$. Thus, by Proposition 4.1.2 (i), $\mu_a << \nu$. 

        Define $\mu_s \equiv \mu - \mu_a$. Thus, $\mu_s(A) = \mu(A \cap (-\infty, 0)$ for all $A \in \borel$. Define $B \equiv (-\infty, 0)$. $\nu(B) = 0$ and $\mu_s(B^c) = \mu_s([0, \infty))= 0$. Thus, $\nu \perp \mu_s$. Thus, $\mu = \mu_a + \mu_s$ is the Lebesgue decomposition of $\mu$ w.r.t. $\nu$.

        $f_{\mu_a} = [N(0, 1)$ pdf] $I_{x \in [0, \infty)}$ \\ $f_\nu = Exp(1)$ pdf

        $\frac{d \mu_a}{d\nu} = \frac{d\mu_a}{d m_1} \left(\frac{d\nu}{d m_1}\right)^{-1} = \frac{f_{\mu_a}}{f_\nu} = \frac{1}{\sqrt{2 \pi}} e^{x - \tfrac{x^2}{2}}I_{x \in [0, \infty)}a.e.(\nu)$
        
       (b) Define $m$ and $m_1$ as in (a). 

       $\mu << m_1 << m$ and from, Example 4.1.1, $m << \nu$. By Proposition 4.1.2 (i), $\mu << \nu$. So, $\mu_a = \mu$. Define $\mu_s \equiv \mu - \mu_a = 0$. $\mu_s$ is the 0 measure, so clearly, $\mu_s \perp \nu$. Therefore the Lebesgue decomposition of $\mu$ w.r.t. $\nu$ is $\mu = \mu_a + \mu_s = \mu$. 

       $f_{\mu_a} = Exp(1)$ pdf \\ $f_\nu = N(0,1)$ pdf.
        $\frac{d \mu_a}{d\nu} = \frac{d\mu_a}{d m} \left(\frac{d\nu}{d m}\right)^{-1} = \frac{f_{\mu_a}}{f_\nu} = \sqrt{2 \pi} e^{x + \tfrac{x^2}{2}}I_{x \in [0, \infty)} a.e.(\nu)$

    (c) From before, we know that $\mu_1 << m$, and since the support of $Cauchy(0, 1)$

    Since $\mu_2$ is discrete and $\nu$ is continuous $\mu_2 \perp \nu$. So, the Lebesgue decomposition of $\mu$ w.r.t. $\nu$ is $\mu = \mu_a + \mu_s$. 

    $f_{\nu} = Cauchy(0,1)$ pdf \\ $f_{\mu_a} = N(0,1)$ pdf.
    
    $\frac{\mu_a}{\nu} =  \frac{d\mu_a}{d m} \left(\frac{d\nu}{d m}\right)^{-1} = \frac{f_{\mu_a}}{f_\nu} = \sqrt{\frac{\pi}{2}}(1 + x^2)e^{x^2/2} a.e.(\nu)$

    (d) 
    Define $\lambda$ as the counting measure 
    $\mu_1 \perp \nu$ because $\mu_1(\nat) = 0$ and $\nu(\nat^c) = 0$
    $\mu_2 << \nu$ because $\mu_2 << \lambda$, $\lambda << \nu$, and this implies $\mu_2 << \nu$

    Thus, $\mu_s = \mu_1$ and $\mu_a = \mu_2$

    $\frac{d \mu_a}{d \nu}(\omega) = \frac{d \mu_a}{d\lambda}(\frac{d \nu}{d\lambda})^{-1} = \frac{\sum_{\omega \in \nat} I_{\omega = k}\frac{1}{k! e}}{\sum_{\omega \in \nat} I_{\omega = k} p(1-p)^k}a.e.(\nu)$

    (e) $\mu_1 << \nu$, because $\mu_1 << m$ and $m << \nu_1$ which implies $m << \nu$ which implies $\mu_1 << \nu$
    Similarly, $\mu_2 << \nu$, because $\mu_2 << \lambda$ and $\lambda << \nu_2$ which implies $\lambda << \nu$ which implies $\mu_1 << \nu$

    By Proposition 4.1.2 (d), $\mu_1 + \mu_2 = \mu << \nu$

    Trivially, $0 \perp \nu$ (where 0 is the zero measure)

    So, $\mu_s = 0$ and $\mu_a = \mu$

     Additionally, it is important to note that $\mu_1 \perp \nu_2$ and $\mu_2 \perp \nu_2$

    From Proposition 4.1.2, for all $A \in \calF$, 
    \begin{align*}
        & \mu(A) = \mu_1(A) + \mu_2(A) \\
        & = \int_A \frac{d\mu_1}{d\nu_1}d\nu_1 + \int_A \frac{d\mu_2}{d\nu_2}d\nu_2 \\
        & = \int_{A\cap \nat^c}  \frac{d\mu_1}{d\nu_1}d\nu_1 + 0 + \int_{A\cap \nat} \frac{d\mu_2}{d\nu_2}d\nu_2 + 0\\
        & = \int_{A\cap \nat^c}  \frac{d\mu_1}{d\nu_1}d(\nu_1 + \nu_2) + \int_{A\cap \nat} \frac{d\mu_2}{d\nu_2}d(\nu_1 + \nu_2)
    \end{align*}

    Thus, $\frac{d \mu_a}{d\nu}(\omega) = \frac{d \mu_1}{d\nu_1}I_{\nat^c}(\omega) +  \frac{d \mu_2}{d\nu_2}I_{\nat}(\omega) = \sqrt{\frac{\pi}{2}}(1 + x^2)e^{-\omega^2/2} I_{\omega \in \nat^c} + \sum_{\omega \in \nat} \frac{1}{e \omega! p(1-p)^\omega}I_{\omega \in \nat}$

    (f) The support of $\mu$ is $\{0, 1, 2, ... , 10\}$ and the support of $\nu$ is $\nat \cap \{0\}$. Clearly, $\mu << \nu$

    Thus, $\mu_a = \mu$ and $\mu_s = 0$, where $\nu = \mu_a + \mu_s$
    For some $A \in \calF$, \\
    $\mu(A) = \sum_{\omega \in A \cap \{0, 1, ... , 10\}} \binom{10}{\omega}(1/2)^{10} = \int_A \frac{d \mu}{d \nu}d\nu = \int_{A \cap \{0, 1, ... , 10\}} \binom{10}{\omega}(1/2)^{10} \frac{1}{e \omega!} d\nu(\omega)$

    Thus, $\frac{d \mu}{d \nu} = \binom{10}{\omega}(1/2)^{10} \frac{1}{e \omega!}I_{\{0, 1, ... , 10\}}(\omega) a.e.(\nu)$
    \end{solution}

%%%%%%%%%%% Question 4 %%%%%%%%%%%%%%%%%%%%    
    \question \textbf{Additional Problem}

    Let $(\Omega, \calF)$ be a measurable space, and $\mu, \: \mu_1, \: \mu_2, \: \nu$ are $\sigma$-finite measures on it. Define the set function $\mu_3 \equiv \mu_1 + \mu_2$ (i.e. for all $A \in \calF$, $\mu_3(A) = \mu_1(A) + \mu_2(A)$. Then prove the following: 
    \begin{enumerate}
        \item[(a)] $\mu_3$ is a measure
        \item[(b)] if $\mu_1 << \nu$ and $\mu_2 << \nu$, then $\mu_3 << \nu$ 
        \item[(c)] Express the Radon-Nikodym derivative $\frac{d \mu_3}{d \nu}$ in terms of $\frac{d \mu_1}{d \nu}$  and $\frac{d \mu_2}{d \nu}$ (proof required).
        \item[(d)] If $\mu << \nu$ and $\mu \perp \nu$ holds, show that $\mu$ is the zero measure.
        \item[(e)] If $\mu_1 = Lognormal(5, 6), \: \mu_2 = Geometric(0.3),$ and $\nu = Poisson(8)$ on $(\Omega, \calF) = (\real , \borel)$, then find the Lebesgue decomposition of $\mu_3$ w.r.t $\nu$, and find the Radon Nikodym derivative $\frac{d \mu_{3a}}{d \nu}$, where $\mu_{3a}$ is the absolutely continuous component of $\mu_3$ in the decomposition.
    \end{enumerate}
    
    (a) \begin{proof}
        (i) Let $A \in \calF$ be given. $\mu_1(A) \in [0, \infty]$ and $\mu_2(A) \in [0, \infty]$, so $\mu_3(A) = \mu_1(A) + \mu_2(A) \in [0, \infty]$

        (ii) $\mu_1(\emptyset) = 0$ and $\mu_2(\emptyset) = 0$, so $\mu_3(\emptyset) = \mu_1(\emptyset) + \mu_2(\emptyset) = 0$

        (iii) Let $A_1, A_2, ... \in \calF$ be disjoint and $\cup_{n \geq 1} A_n \in \calF$. \\
        $\mu_3(\cup_{n \geq 1} A_n) = \mu_1(\cup_{n \geq 1} A_n) + \mu_2(\cup_{n \geq 1} A_n) = \cutesum{n} \mu_1(A_n) + \cutesum{n} \mu_2(A_n) = \cutesum{n} (\mu_1(A_n) + \mu_2(A_n)) = \cutesum{n} \mu_3(A)$

        Thus, $\mu_3$ is a measure on $(\calF, \Omega)$
    \end{proof}

    (b) \begin{proof}
        Let $A \in \calF$ be an arbitrary set such that $\nu(A) = 0$. Since $\mu_1 << \nu$, $\mu_1(A) = 0$. Similarly, $\mu_2(A) = 0$. Thus, $\mu_3(A) = \mu_1(A) + \mu_2(A) = 0 + 0 = 0$. 

        Thus, $\nu(A) = 0 \: \implies \: \mu_3(A) = 0$, and thus $\mu_3 << \nu$
        \end{proof}

    (c) \begin{proof}
        The R-N derivative of $\mu_1$ w.r.t. $\nu$, denoted by $\frac{d \mu_1}{d\nu}$ is a nonnegative function such that $\mu_1(A) = \int_A \frac{d \mu_1}{d\nu} d\nu \: \forall \: A \in \calF$. \\
        The R-N derivative of $\mu_2$ w.r.t. $\nu$, denoted by $\frac{d \mu_2}{d\nu}$ is a nonnegative function such that $\mu_2(A) = \int_A \frac{d \mu_2}{d\nu} d\nu \: \forall \: A \in \calF$. \\
        The R-N derivative of $\mu_3$ w.r.t. $\nu$, denoted by $\frac{d \mu_3}{d\nu}$ is a nonnegative function such that $\mu_3(A) = \int_A \frac{d \mu_3}{d\nu} d\nu \: \forall \: A \in \calF$.

        For all $A \in \calF$, $\mu_3(A) = \mu_1(A) + \mu_2(A) = \int_A \frac{d \mu_1}{d\nu} d\nu + \int_A \frac{d \mu_2}{d\nu} d\nu \\ = \int_A \frac{d \mu_1}{d\nu} + \frac{d \mu_2}{d\nu} d\nu  = \int_A \frac{d \mu_3}{d\nu} d\nu$

        Thus, $\frac{d \mu_3}{d\nu} = \frac{d \mu_1}{d\nu} + \frac{d \mu_2}{d\nu}$.
        \end{proof}

        (d) \begin{proof}
            Since $\mu \perp \nu$, there must exist some $B \in \calF$ such that $\nu(B) = 0$ and $\mu(B^c) = 0$. Since $\mu << \nu$, this means that $\mu(B) = 0$.

            $B, B^c \in \calF$ and $B \cap B^c = \empty$, so, by the def of a measure, $\mu(B) + \mu(B^c) = \mu(B \cup B^c) = \mu(\Omega) = 0$.

            $A \subset \Omega$ for all $A \in \calF$, so by monotonicity, $\mu(A) \leq \mu(\Omega) = 0$. Since $\mu(A) \in [0, \infty]$, $\mu(A) = 0$ for all $A \in \calF$. Thus, $\mu$ is the zero measure. 
        \end{proof}

\newpage

        (e)
        \begingroup

\setlength{\tabcolsep}{10pt} % Default value: 6pt
\renewcommand{\arraystretch}{1.5} % Default value: 1
        \begin{table}[h!]
            \centering
            \begin{tabular}{|c|c|c|c|}
            \hline
            \textbf{Measure} & \textbf{Support} & $g(x)$ & $f(x)$ \\
            \hline
            $\mu_1$ &  $(0, \infty)$ & $g_{\mu_1}(x) = \frac{1}{x \sqrt{12 \pi}}\exp\left(-\frac{(\log(x) - 6)^2}{12}\right)$ & $f_{\mu_1} = g_{\mu_1}(x)I_{x \in (0, \infty)}$\\
            \hline
            $\mu_2$ & $\integer_+$ & $g_{\mu_2}(x) = 1 - (0.7)^{x +1}$ & $f_{\mu_2}(x) = g_{\mu_2}(x)I_{x \in \integer_+}$\\
            \hline
            $\nu$ & $\integer_+$ & $g_{\nu}(x) = \frac{8^x e^{-8}}{x!}$ & $f_{\nu}(x) = g_{\nu}(x)I_{x \in \integer_+}$\\
            \hline
            \end{tabular}
        \end{table}
        \endgroup

        $\lambda =$ the counting measure

        If we define 

        $\mu_2 << \nu$ and $\mu_1 \perp \nu$, so $\mu_{3a} = \mu_2$. 
         $\mu_2 << \lambda << \nu$, thus, by Proposition 4.1.2, $\frac{d \mu_{3a}}{d \nu} = \frac{d\mu_{3a}}{d\lambda}\frac{d\lambda}{d \nu} = \frac{d\mu_{3a}}{d\lambda}\frac{d\nu}{d \lambda})^{-1} = \frac{f_{\mu_2}}{f_\nu}$

\end{questions}
\end{document}